%Este trabalho está licenciado sob a Licença Creative Commons Atribuição-CompartilhaIgual 3.0 Não Adaptada. Para ver uma cópia desta licença, visite https://creativecommons.org/licenses/by-sa/3.0/ ou envie uma carta para Creative Commons, PO Box 1866, Mountain View, CA 94042, USA.

\chapter{Introdução}\label{chap:introducao}
\emconstrucao

\section{Números reais e intervalos de numéricos reais}\label{sec:intro_reais}
\construirSec

\subsection*{Exercícios resolvidos}

\construirExeresol


\subsection*{Exercícios}

\construirExer


%***********************************************************************************
\section{Funções reais e gráfico}\label{sec:intro_fun_reais}
\construirSec

\subsection*{Exercícios resolvidos}

\construirExeresol


\subsection*{Exercícios}

\construirExer


%***********************************************************************************
\section{Funções trigonométricas}\label{sec:intro_trigo}
\construirSec

\subsection*{Exercícios resolvidos}

\construirExeresol


\subsection*{Exercícios}

\construirExer


%***********************************************************************************
\section{Função inversa}\label{sec:intro_inversa}
\construirSec

Sobre funções inversas é importante conhecer algumas definições, como por exemplo:
\frame{
1. Definição: Uma função $f$ que nunca assume o mesmo valor duas vezes é chamada de função injetora, ou seja,
\begin{center}
$f(x_{1}) \neq f(x_{2})$ se $x_{1} \neq x_{2}$
\end{center}
}

\frame{
2. Definição: Seja a função $f$ uma função injetora com domínio $A$ e imagem $B$. Logo, a sua função inversa $f^{-1}$ tem domínio $B$ e imagem $A$, sendo definida por
\begin{center}
$f^{-1}(y) = x \Leftrightarrow f(x) = y \forall y$ em $B$
\end{center}
}

\subsection*{Exercícios resolvidos}

\construirExeresol


\subsection*{Exercícios}

\construirExer


%***********************************************************************************
\section{Função composta}\label{sec:intro_composta}
\construirSec


\subsection*{Exercícios resolvidos}

\construirExeresol

\subsection*{Exercícios}

\construirExer

\section{Exercícios finais}

\construirExer
